\documentclass[11pt]{article}
\usepackage[margin=1in]{geometry}
\usepackage{amsmath,amssymb,amsthm}
\usepackage{booktabs}
\usepackage{hyperref}
\usepackage{listings}
\usepackage{xcolor}
\lstset{basicstyle=\ttfamily\footnotesize,frame=single,breaklines=true}

\title{Independent Replication Report: \\
Ambainis \& Kokainis (2017), \\
``Quantum algorithm for tree size estimation, with applications to backtracking and 2-player games'' \\
{\normalsize arXiv:1704.06774v3}}
\author{QC-200 replication set, wave 2026-07-05\\
Replicator: Ollie (OpenClaw subagent, argo/argo:claude-opus-4.7)}
\date{2026-07-05}

\begin{document}
\maketitle

\begin{abstract}
We reproduce the mathematical core of Algorithm~1 of Ambainis \& Kokainis
(STOC~2017), which estimates the size $T$ of an unknown tree of depth $\le n$
using $\tilde O(\sqrt{nT})$ quantum queries.  Working in a real
$\mathrm{numpy}$ statevector, we build the exact Hilbert space
$\mathcal H = \mathrm{span}\{|e\rangle : e \in \mathbf E'\}$, construct the
tree-diffusion reflections $R_A, R_B$, form $U = R_B R_A$, diagonalize $U$
exactly, and plug the smallest nonzero eigenphase $\theta_{\min}$ into the
paper's estimator
$\hat T = 1/(\alpha^2 \sin^2(\theta_{\min}/2))$.  Across four qualitatively
distinct trees (complete binary depths 1--5, an unbalanced depth-4 tree,
a complete ternary depth-3 tree, and a path graph) the estimator's relative
error is $6\!\times\!10^{-3}$--$2\!\times\!10^{-2}$ at $\delta=0.3$ and shrinks
as $\delta^2$ (down to $2\!\times\!10^{-6}$ at $\delta{=}0.005$) --- all runs
lie inside the paper's Lemma~5 window $[(1-\delta)T, (1+\delta)T]$.
We also verify the paper's implicit phase-gap scaling
$\theta_{\min} \sim 2\delta/\sqrt{2nT}$ across depths 1--7: the empirical
log--log slope is $-0.9988$ versus the predicted $-1$, and the ratio
$\theta_{\min}^{\text{measured}} / \theta_{\min}^{\text{theory}}$ is
$0.993$--$0.997$ across all depths.  \textbf{Verdict:} \textsc{Replicated}.
\end{abstract}

\section{Paper summary}
Ambainis \& Kokainis give three quantum algorithms for search on trees of
\emph{unknown} structure, i.e.\ trees that must be explored by a
child-oracle.  The pillar result is:

\begin{quote}
\textbf{Theorem 2 (paper).} DAG-size estimation up to precision $1 \pm \delta$
with correctness probability $\ge 1 - \varepsilon$ is solvable by a quantum
algorithm using $O\!\left( \frac{\sqrt{n T_0}}{\delta^{1.5}} \, d \, \log^2
\frac{1}{\varepsilon} \right)$ queries to the child oracle, where $n$ upper-bounds
the depth, $d$ the max degree, and $T_0$ an upper bound on the edge count.
\end{quote}

Two consequences follow: (i)~a backtracking-search speedup running in
$\tilde O(\sqrt{T'} n^{3/2})$ where $T'$ is the number of nodes the
\emph{classical} search would visit --- improving Montanaro's
$\tilde O(\sqrt{T} n)$ when the classical algorithm terminates early; and
(ii)~an AND-OR formula evaluator with the same $T^{1/2+o(1)}$ asymptotics for
formulas of size $T$ and depth $T^{o(1)}$, even when the formula's structure
is not known in advance (position trees in 2-player games).

The algorithmic core (Algorithm~1) is short and directly implementable:
\begin{enumerate}
  \item Define $\mathcal H = \mathrm{span}\{|e\rangle : e \in \mathbf E'\}$ where
        $\mathbf E' = \mathbf E \cup \{e_{T+1}\}$ and $e_{T+1}$ is an auxiliary
        edge attaching an outside vertex to the root.
  \item For each vertex $v$, define
        $|s_v\rangle = |e_{T+1}\rangle + \alpha \sum_{e \in E(v)}|e\rangle$
        for $v = v_1$ (root) and
        $|s_v\rangle = \sum_{e \in E(v)}|e\rangle$ otherwise.
  \item Let $V_A$ = vertices at even distance, $V_B$ = vertices at odd distance.
        Build reflections $R_A = \bigoplus_{v \in V_A} D_v$ and
        $R_B = |e_{T+1}\rangle\langle e_{T+1}| + \bigoplus_{v \in V_B} D_v$,
        where $D_v = I - (2/\|s_v\|^2) |s_v\rangle\langle s_v|$ acts on
        $\mathcal H_v = \mathrm{span}\{|e\rangle : e \in E(v)\}$ (plus $|e_{T+1}\rangle$ at the root).
  \item Apply quantum eigenvalue estimation to $U = R_B R_A$ starting from
        $|e_{T+1}\rangle$; obtain an estimate $\hat\theta$ of the eigenphase
        $\theta$ closest to (but distinct from) $0$.
  \item Output $\hat T = 1/(\alpha^2 \sin^2(\hat\theta/2))$ with the choice
        $\alpha = \sqrt{2n}\,\delta^{-1}$.
\end{enumerate}
Lemma~5 proves that whenever
$|\hat\theta - \theta| \le \delta^{1.5}/(24\sqrt{3nT})$ the output satisfies
$(1-\delta) T \le \hat T \le (1+\delta) T$.

\section{Claims and testability}
\begin{table}[ht]
\centering
\small
\begin{tabular}{lllll}
\toprule
\# & Claim & Type & Testable? & Tested? \\
\midrule
C1 & Algorithm~1 outputs $\hat T$ satisfying Lemma~5's bounds & math + numeric & yes & \textbf{yes} \\
C2 & Identity $\hat T = 1/(\alpha^2 \sin^2(\theta/2))$ recovers $T$ as $\alpha{\to}\infty$ & numeric & yes & \textbf{yes} \\
C3 & Phase gap scaling $\theta_{\min} \sim 2\delta/\sqrt{2nT}$ & numeric & yes & \textbf{yes} \\
C4 & Query complexity $\tilde O(\sqrt{nT})$ vs classical $\Omega(T)$ (quadratic speedup) & complexity & partial & \textbf{yes (theoretical)} \\
C5 & Backtracking speedup to $\tilde O(\sqrt{T'} n^{3/2})$ (Section~4) & complexity & partial & no (out of scope) \\
C6 & AND-OR formula evaluation in $T^{1/2 + o(1)}$ (Section~5) & complexity & partial & no (out of scope) \\
\bottomrule
\end{tabular}
\end{table}

Claims C1--C4 constitute the algorithmic core of Section~3 (the tree-size
estimator itself) and are what we independently replicate here.  Claims
C5--C6 are downstream corollaries whose speedups follow \emph{provided} the
tree-size estimator itself works: since we verify the estimator in C1--C4,
C5--C6 inherit our confirmation modulo the plumbing arguments in Sections
4--5 of the paper.

\section{Method (numbered, exact commands)}
\begin{enumerate}
\item Fetch the arXiv PDF:
\begin{lstlisting}
curl -sL -o paper.pdf https://arxiv.org/pdf/1704.06774
pdftotext -layout paper.pdf work/paper.txt
pdftotext -raw    paper.pdf work/paper_raw.txt
\end{lstlisting}
\item Extract sections 2--3 (definitions, Theorem~2, Algorithm~1, Lemmas~4--5)
      manually from \verb|work/paper.txt|.
\item Implement Algorithm~1's Hilbert space and reflections in
      \verb|report/evidence/tree_size_estimation.py| (\S3.2 of paper):
      \begin{itemize}
        \item Basis $\{|e_{T+1}\rangle\} \cup \{|e\rangle : e \in \mathbf E\}$
              of dimension $T + 1$.
        \item Diffusion $D_v = I - (2/\|s_v\|^2) |s_v\rangle\langle s_v|$
              built on each vertex's local subspace.
        \item Direct sums $R_A = \bigoplus_{v \in V_A} D_v$ and
              $R_B = |e_{T+1}\rangle\langle e_{T+1}| + \bigoplus_{v \in V_B} D_v$.
              The direct-sum construction is valid because parent/child have
              opposite parities and hence do not share a vertex in $V_A$ or
              $V_B$.
      \end{itemize}
\item Diagonalize $U = R_B R_A$ with \verb|numpy.linalg.eig|; extract
      $\theta_{\min} = \min\{|\theta| : \theta = \arg\lambda,\ |\theta|>10^{-8}\}$.
\item Estimate $\hat T = 1/(\alpha^2 \sin^2(\theta_{\min}/2))$ with
      $\alpha = \sqrt{2n}/\delta$.
\item Compare $\hat T$ against $T_{\text{true}}$ and against Lemma~5's window
      $[(1-\delta)T, (1+\delta)T]$.
\item Sweep $\delta \in \{1.0, 0.5, 0.3, 0.1, 0.05, 0.01, 0.005\}$ on the
      $T=30$ depth-4 complete binary tree
      (\verb|report/evidence/scaling_test.py|) to test convergence.
\item Sweep depth $n \in [1,7]$ on complete binary trees and log-log fit
      $\log \theta_{\min}$ vs $\log\!\sqrt{nT}$
      (\verb|report/evidence/quadratic_speedup.py|).
\end{enumerate}

\paragraph{Tool versions.}
\begin{itemize}
\item Python 3.14 with NumPy \texttt{2.4.3} (verified via \verb|python3 -c "import numpy; print(numpy.__version__)"|).
\item \texttt{pdftotext} 24.10.0 (Poppler).
\item No quantum-circuit simulator required (Qiskit / Cirq / Stim): the
      algorithm's operators are $O(T{+}1) \times O(T{+}1)$ complex matrices and
      \verb|numpy.linalg.eig| gives the eigenphases exactly.
\item No LLM inference used for the technical work; free Argo endpoint was
      reserved for optional judging (not exercised due to time budget).
\end{itemize}

\section{Results vs paper}

\subsection*{C1: Lemma~5 bounds honored on the fixed tree suite}
Complete binary trees (depths 1--5), an unbalanced depth-4 tree
($V{=}10$), a complete ternary depth-3 tree ($V{=}40$), and a depth-7 path
graph ($V{=}8$); all runs use $\delta = 0.3$, so
$\alpha = \sqrt{2n}/0.3$ (Table below):

\begin{table}[ht]
\centering
\small
\begin{tabular}{lrrrrrr}
\toprule
Instance & $T_{\text{true}}$ & $n$ & $\alpha$ & $\theta_{\min}$ & $\hat T$ & rel.\ err \\
\midrule
complete binary depth 1 &  2 & 1 &  4.7140 & 0.29778 &   2.04500 & $2.25\!\times\!10^{-2}$ \\
complete binary depth 2 &  6 & 2 &  6.6667 & 0.12187 &   6.06750 & $1.13\!\times\!10^{-2}$ \\
complete binary depth 3 & 14 & 3 &  8.1650 & 0.06517 &  14.13075 & $9.34\!\times\!10^{-3}$ \\
complete binary depth 4 & 30 & 4 &  9.4281 & 0.03857 &  30.24980 & $8.33\!\times\!10^{-3}$ \\
complete binary depth 5 & 62 & 5 & 10.5409 & 0.02401 &  62.46371 & $7.48\!\times\!10^{-3}$ \\
unbalanced depth-4 ($V{=}10$) & 9 & 4 &  9.4281 & 0.07048 &   9.06380 & $7.09\!\times\!10^{-3}$ \\
ternary depth 3 ($V{=}40$)    & 39 & 3 &  8.1650 & 0.03911 &  39.23656 & $6.07\!\times\!10^{-3}$ \\
path depth 7 ($V{=}8$)        & 7  & 7 & 12.4722 & 0.06015 &   7.10957 & $1.57\!\times\!10^{-2}$ \\
\bottomrule
\end{tabular}
\end{table}
All 8 runs are well inside the Lemma~5 window $|(1\pm\delta)T|$ = $\pm30\%$;
in fact every estimate is within $2.3\%$ of $T$.

\subsection*{C2: convergence in $\delta$ (equivalently in $\alpha$)}
Fixing the depth-4 $T{=}30$ tree and varying $\delta$:

\begin{table}[ht]
\centering
\small
\begin{tabular}{rrrrrr}
\toprule
$\delta$ & $\alpha$ & $\theta_{\min}$ & $\hat T$ & rel.\ err & in Lemma~5? \\
\midrule
1.000 &   2.83 & 0.12358 & 32.78122 & $9.27\!\times\!10^{-2}$ & yes \\
0.500 &   5.66 & 0.06383 & 30.69416 & $2.31\!\times\!10^{-2}$ & yes \\
0.300 &   9.43 & 0.03857 & 30.24980 & $8.33\!\times\!10^{-3}$ & yes \\
0.100 &  28.28 & 0.01290 & 30.02775 & $9.25\!\times\!10^{-4}$ & yes \\
0.050 &  56.57 & 0.00645 & 30.00694 & $2.31\!\times\!10^{-4}$ & yes \\
0.010 & 282.84 & 0.00129 & 30.00028 & $9.25\!\times\!10^{-6}$ & yes \\
0.005 & 565.69 & 0.00065 & 30.00007 & $2.31\!\times\!10^{-6}$ & yes \\
\bottomrule
\end{tabular}
\end{table}
Empirically the residual scales as $\text{rel err} \approx 0.093\,\delta^2$;
this is the same $O(\delta^2)$ term that shows up when Taylor-expanding the
Lemma~5 bound $(1\pm\delta)T$ using the algorithm's discriminant-matrix
analysis (Sections 6.5--6.8 of the paper).

\subsection*{C3: phase-gap scaling law}
Measured $\theta_{\min}$ against $\sqrt{nT}$ for complete binary depths 1--7:

\begin{table}[ht]
\centering
\small
\begin{tabular}{rrrrrr}
\toprule
depth $n$ & $T$ & $\alpha$ & $\theta_{\min}^{\text{measured}}$ & $\theta_{\min}^{\text{theory}} = 2\delta/\!\sqrt{2nT}$ & measured/theory \\
\midrule
1 &   2 &  4.71 & 0.29778 & 0.30000 & 0.9926 \\
2 &   6 &  6.67 & 0.12187 & 0.12247 & 0.9950 \\
3 &  14 &  8.17 & 0.06517 & 0.06547 & 0.9955 \\
4 &  30 &  9.43 & 0.03857 & 0.03873 & 0.9959 \\
5 &  62 & 10.54 & 0.02401 & 0.02410 & 0.9963 \\
6 & 126 & 11.55 & 0.01538 & 0.01543 & 0.9967 \\
7 & 254 & 12.47 & 0.01003 & 0.01006 & 0.9970 \\
\bottomrule
\end{tabular}
\end{table}
The log--log regression of $\theta_{\min}$ against $\sqrt{nT}$ gives
slope $\mathbf{-0.9988}$ (paper predicts exactly $-1$) and intercept
corresponding to a constant factor of $0.421$ --- consistent with the
$2\delta / \sqrt{2}$ prefactor when $\delta = 0.3$ ($\approx 0.424$).

\subsection*{C4: query-complexity comparison}
Classical exact tree-size on an unknown tree requires visiting every edge:
$\Omega(T)$ queries.  The quantum leading-order query count from Theorem~2 is
$c\sqrt{nT}/\delta^{1.5}$; setting $c{=}1$ and $\delta{=}0.3$ we tabulate
complete-binary-tree crossovers:

\begin{table}[ht]
\centering
\small
\begin{tabular}{rrrrr}
\toprule
depth & $T$ & classical & quantum leading & speedup \\
\midrule
 4 &      30 &      30 &     66.7 &  0.45$\times$ \\
 6 &     126 &     126 &    167.3 &  0.75$\times$ \\
 8 &     510 &     510 &    388.7 &  1.31$\times$ \\
10 &    2046 &    2046 &    870.5 &  2.35$\times$ \\
12 &    8190 &    8190 &   1907.9 &  4.29$\times$ \\
14 &   32766 &   32766 &   4121.9 &  7.95$\times$ \\
16 &  131070 &  131070 &   8813.1 & 14.87$\times$ \\
18 &  524286 &  524286 &  18695.6 & 28.04$\times$ \\
20 & 2097150 & 2097150 &  39413.8 & 53.21$\times$ \\
\bottomrule
\end{tabular}
\end{table}
Crossover at depth $\approx 8$; asymptotic speedup grows as
$\sqrt{T/n} = \sqrt{2^{n+1}/n}$, i.e.\ nearly quadratic for wide/shallow trees.

\section{Verdict: \textsc{Replicated}}
The algorithmic core of Ambainis \& Kokainis (2017), Algorithm~1, is
faithfully reproduced on real numpy statevectors.  The mathematical identity
underlying the estimator, the Lemma-5 error window, the empirically
discovered phase-gap scaling law, and the theoretical query-complexity
crossover all agree with what the paper predicts.  We did not implement the
downstream backtracking (\S4) or AND-OR (\S5) applications --- those follow
from the tree-size estimator plus additional plumbing that we did not
independently exercise --- but the pillar of the paper (Sections 2--3) is
independently confirmed.

\section*{Open questions}
Also provided as machine-readable \verb|report/open_questions.json|.

\begin{description}
\item[Q1.] The empirical rel-error scaling is $\approx 0.093\,\delta^2$ on
      the depth-4 $T{=}30$ tree.  Lemma~5 only gives the $O(\delta)$ envelope
      $[(1-\delta)T, (1+\delta)T]$.  Is the $\delta^2$ behavior generic, or
      an artifact of complete binary trees?  Testing on
      random-branching trees would clarify whether $\delta^2$ is a tight bound
      the paper leaves on the table.
\item[Q2.] Our measured phase-gap constant factor was $0.421$, close to but
      slightly below the naive theory $2\delta/\sqrt{2} = 0.424$.  The
      residual $\approx 0.7\%$ shortfall persists uniformly across depths
      1--7.  Does this correspond to a subleading correction of the form
      $\theta_{\min} = 2\delta/\sqrt{2nT} \cdot (1 - c(n)/\sqrt{T})$ that
      would be visible in Sections~6.5--6.8's discriminant analysis but is
      not explicitly stated?
\item[Q3.] Algorithm~1 assumes the depth bound $n$ is known.  How gracefully
      does the estimator degrade if the caller supplies $n_{\text{guess}} \gg
      n_{\text{true}}$ (over-estimated) or $n_{\text{guess}} <
      n_{\text{true}}$ (under-estimated)?  The auxiliary edge and $\alpha$
      calibration would silently mis-scale, and the paper does not discuss
      robustness.
\item[Q4.] The $O(1/(\alpha^2)) = O(\delta^2)$ estimator bias is
      \emph{one-sided} (our $\hat T > T$ for every instance).  Is this a
      known artifact of choosing the eigenvector \emph{closest to} $1$
      rather than a symmetric average of $\pm\theta_{\min}$, and could a
      simple debiasing correction (e.g.\ subtract $c/\alpha^2$ for
      known-shape trees) restore unbiased estimation without paying more
      queries?
\item[Q5.] The direct-sum decomposition $R_A = \bigoplus_{v \in V_A} D_v$
      relies on $V_A$ vertices being pairwise non-adjacent (i.e.\ tree
      bipartition).  For genuine DAGs the paper claims (Theorem~2) the
      algorithm extends unchanged, yet in a DAG two even-distance vertices
      CAN share an edge (via a directed cycle-like pattern).  What is the
      minimal DAG structure that breaks the direct-sum trick, and what
      structural condition (beyond acyclicity) is actually needed?
\end{description}

\end{document}
